\begingroup
\setlength{\parindent}{0pt}
\setlength{\parskip}{0.52\baselineskip}
\setlength{\abovedisplayskip}{8.5pt plus 2pt minus 2pt}
\setlength{\belowdisplayskip}{8.5pt plus 2pt minus 2pt}
\setlength{\jot}{6.5pt}

\section*{How the Proof Works}

One smooth three-dimensional incompressible flow can contain a vortex, stretch
it, and drive its core toward finer scales.  The global danger is whether that
flow can sharpen one of its own structures so far that smooth continuation
fails in finite time.  Fix one velocity \(u\), one pressure \(p\), and one
viscosity \(\nu>0\), and follow only that history:
\[
 \partial_tu+(u\!\cdot\nabla)u+\nabla p=\nu\Delta u,
 \qquad \nabla\!\cdot u=0,
 \qquad u(0)=u_0.
\]
The transport term moves and deforms the velocity already in the fluid.
Pressure responds across the whole field so that this motion remains
incompressible.  Viscosity diffuses a sharp change into its surroundings.

Take one vortex tube inside that field.  When the surrounding strain pulls the
tube along its axis, incompressibility forces a transverse contraction: a
longer piece of fluid cannot keep its old cross-section without gaining
volume.  The core narrows.  At the level of vorticity
\(\omega=\nabla\times u\), this motion is visible in
\[
 \partial_t\omega+(u\!\cdot\nabla)\omega
 = (\omega\!\cdot\nabla)u+\nu\Delta\omega.
\]
The stretching term can amplify rotation aligned with the strain.  As the
core contracts, the velocity turns across a shorter transverse distance and
its gradient steepens.  Nonlinear motion can sharpen the vortex in this way,
while \(\nu\Delta\omega\) spreads the concentrated vorticity and smooths that
gradient.  The narrower the core becomes, the stronger both effects can
become.

A narrowing vortex also changes the length and time scales on which the
equation is being read.  If \((u,p)\) is viewed a factor \(\lambda\) deeper in
space, then
\[
 u_\lambda(x,t)=\lambda u(\lambda x,\lambda^2t),
 \qquad
 p_\lambda(x,t)=\lambda^2p(\lambda x,\lambda^2t).
\]
Space has contracted by \(\lambda^{-1}\), the clock by \(\lambda^{-2}\), and
the velocity amplitude has grown by \(\lambda\).  Every term in the equation
then carries a factor \(\lambda^3\).  This is what Navier--Stokes scaling says:
it compares solutions seen at different resolutions.  It does not say that an
evolving vortex must pass through a chain of rescaled copies.

The spatial Sobolev norms record what that comparison does to regularity:
\[
 \|u_\lambda(\cdot,t)\|_{\dot H^s}
 =\lambda^{s-1/2}
  \|u(\cdot,\lambda^2t)\|_{\dot H^s}.
\]
Below \(s=\tfrac12\), finer concentration lowers the scaled norm.  Above
\(s=\tfrac12\), it raises it.  At \(s=\tfrac12\), it leaves the scaled size
unchanged.  Thus \(\dot H^{1/2}\) is scaling-critical; it is not a conserved
quantity.  Saying that nonlinear transport and viscosity have equal scaling
weight means only that rescaling cannot make one term disappear beside the
other.  It gives no cancellation, sign, or bound.  Kinetic energy behaves
differently:
\[
 \|u_\lambda(\cdot,t)\|_2^2
 =\lambda^{-1}\|u(\cdot,\lambda^2t)\|_2^2.
\]
The idealized core can move to finer scales while its energy remains finite
and its higher gradients grow.

Now suppose this sharpening repeats by a fixed spatial factor \(\lambda>1\).
After \(j\) contractions, the scaling above gives a core length and a
parabolic response time
\[
 \ell_j=\lambda^{-j}\ell_0,
 \qquad
 \tau_j=\lambda^{-2j}\tau_0,
 \qquad
 \sum_{j=0}^{\infty}\tau_j
 =\frac{\tau_0}{1-\lambda^{-2}}<\infty.
\]
Infinitely many scale changes can fit inside a finite \(T_*\).  The resulting
Zeno cascade is a picture of one finite-energy fluid reaching arbitrarily
small scales while its gradients become arbitrarily large, not a theorem that
a vortex actually follows this path.

Inside that hypothetical velocity-blow-up picture, let \(X\) be the uniform
norm in which \(u\) is supposed to diverge.  If its acceleration stayed
bounded in \(X\), then the fundamental theorem of calculus would give
\[
 \sup_{t<T_*}\|u(t)\|_X
 \le \|u(0)\|_X
   +T_*\sup_{t<T_*}\|\partial_tu(t)\|_X<\infty,
\]
so the pictured velocity blow-up could not occur.  Applied one order higher,
bounded jerk keeps acceleration bounded; the calculation repeats at every
finite order.  In this picture, in order for a singularity to happen, velocity
has to blow up.  In order for that to happen, its acceleration would also have
to blow up, which means its jerk has to blow up, and so on, and so on:
\[
 u,\qquad \partial_tu,\qquad \partial_t^2u,\qquad\ldots
\]
The hypothetical cascade cannot stop at any fixed temporal order.

The equation carries this repeated temporal climb into its nonlinear and
viscous structure:
\[
\begin{aligned}
 \partial_tu
 &=\nu\Delta u-(u\!\cdot\nabla)u-\nabla p,\\
 \partial_t^2u
 &=\nu\Delta\partial_tu
   -(\partial_tu\!\cdot\nabla)u
   -(u\!\cdot\nabla)\partial_tu
   -\nabla\partial_tp,\\
 \partial_t^3u
 &=\nu\Delta\partial_t^2u
   -(\partial_t^2u\!\cdot\nabla)u
   -2(\partial_tu\!\cdot\nabla)\partial_tu
   -(u\!\cdot\nabla)\partial_t^2u
   -\nabla\partial_t^2p.
\end{aligned}
\]
The Laplacian follows the rising temporal response, while the nonlinear term
opens into every product allowed by the product rule.  After three lines the
bookkeeping is already hiding the relation, so put
\[
 V_r:=\partial_t^ru,
 \qquad
 Q_r:=\partial_t^rp.
\]
At order \(r\), those product-rule copies occur with multiplicity
\(\binom ra\), so the expansion closes as
\[
\begin{aligned}
 V_{r+1}
 &+\sum_{a=0}^{r}\binom ra
   (V_a\!\cdot\nabla)V_{r-a}
   +\nabla Q_r
   =\nu\Delta V_r,\\
 Q_r
 &=R_iR_j\sum_{a=0}^{r}\binom ra
   V_{a,i}V_{r-a,j}.
\end{aligned}
\]
Incompressibility fixes \(Q_r\) through the Riesz transforms \(R_i\), so the
pressure row stays with the nonlinear row at every order.

If a singularity were forming, one would expect spatial regularity to erode as
these temporal responses rise.  Yet each next response contains
\(\nu\Delta V_r\), hence two spatial derivatives of the preceding response.
That does not give regularity on its own.  It tells us which spatial
coordinates the equation exposes.

To measure the spatial order exposed by \(\Delta V_r\), let
\(\Lambda=(-\Delta)^{1/2}\) and set
\[
 E_s(t):=\frac12\|\Lambda^su(t)\|_2^2,
 \qquad
 H(t):=E_{1/2}(t),
\]
and retain the complete pressure-coupled nonlinear transfer
\[
 \mathcal P_s(t)
 :=-\operatorname{Re}\left\langle
 \Lambda^su,
 \Lambda^s\bigl((u\!\cdot\nabla)u+\nabla p\bigr)
 \right\rangle_{L^2}.
\]
Pairing the equation at height \(s\) gives
\[
 E_s'=\mathcal P_s-2\nu E_{s+1}.
\]
Differentiate this identity and use it again at each new height.  At every
finite order \(n\),
\[
 H^{(n)}
 =(-2\nu)^nE_{n+1/2}
 +\sum_{j=0}^{n-1}(-2\nu)^j
   \partial_t^{\,n-1-j}\mathcal P_{j+1/2}.
\]
Acceleration, jerk, and the higher temporal responses have now exposed the
Sobolev coordinates
\[
 E_{n+1/2}
 =\frac12\|u\|_{\dot H^{n+1/2}}^2.
\]
If every finite exposed height reaches one terminal value \(u_*\), with the
finite-energy row retained, then for any \(k\) one can choose \(n>k+1\) and
read
\[
 u_*\in L^2\cap\bigcap_{n<\infty}\dot H^{n+1/2}
 =H^\infty
 \hookrightarrow C^\infty.
\]
But exposure is not control.  Differentiation does not bound
\(H^{(n)}\), the transfer terms, or \(E_{n+1/2}\); it only keeps their exact
relation visible.  The singular picture predicts loss of spatial regularity
as temporal order rises, and these identities alone do not prevent it.

Let \([0,T_*)\) be the maximal classical lifespan and suppose
\(T_*<\infty\).  Fix one finite pair \((N,M)\) and one cutoff
\(0<T<T_*\).  For \(0\le n\le N\), keep the whole transfer
\[
 \mathcal B_n
 :=\sum_{a=0}^{n}\binom na(V_a\!\cdot\nabla)V_{n-a}+\nabla Q_n.
\]
Pairing the recurrence for \(V_n\) at height \(M\) gives
\[
\begin{aligned}
 \frac12\frac{d}{dt}\|\Lambda^M V_n\|_2^2
 &+\operatorname{Re}\left\langle
   \Lambda^M V_n,\Lambda^M\mathcal B_n
  \right\rangle_{L^2}\\
 &=\nu\operatorname{Re}\left\langle
   \Lambda^M V_n,\Lambda^M\Delta V_n
  \right\rangle_{L^2}
 =-\nu\|V_n\|_{\dot H^{M+1}}^2.
\end{aligned}
\]
The viscous term has a fixed sign, while every product-rule transfer and the
pressure response remain coupled to it in this finite system.  This identity
shows what must be controlled; it does not supply the whole-terminal bound by
itself.  That payment is the datum-generated theorem:
\[
 \sup_{0<T<T_*}\sum_{n=0}^{N}\left(
  \|V_n\|_{L^2(0,T;H^{M+1})}^2
  +\|V_{n+1}\|_{L^2(0,T;H^{M-1})}^2
 \right)<\infty.
\]
Every term in the sum increases monotonically with the cutoff.  Since one
finite bound holds for all \(T<T_*\), the limit \(T\uparrow T_*\) places each
adjacent pair on the full approach to \(T_*\).  With
\(V_{n+1}=\partial_tV_n\), that pair sits in the Hilbert triple
\[
 H^{M+1}\subset H^M\subset H^{M-1}
\]
and gives
\[
 V_n\in C([0,T_*];H^M).
\]
The pressure formula supplies
\(Q_n\in L^1(0,T_*;H^M)\) and
\(\nabla Q_n\in L^1(0,T_*;H^{M-1})\), so the finitely many rows now form one
pressure-complete rectangle
\[
 \mathcal R_{N,M}[u_0;T_*]
 =\bigl((V_n)_{n=0}^{N},(V_{n+1})_{n=0}^{N};(Q_n)_{n=0}^{N}\bigr).
\]
The bound may depend on \(N\) and \(M\).  No constant uniform over every
derivative order is being used.

A larger rectangle retains each older coordinate as the identical
distributional derivative of \(u\) or \(p\).  Restricting from
\((N,M)\) to \((N',M')\) consequently gives
\[
 \rho_{N',M'}^{N,M}\mathcal R_{N,M}[u_0;T_*]
 =\mathcal R_{N',M'}[u_0;T_*].
\]
All finite rectangles can now meet without changing the fluid on an overlap.
The finite coordinate space containing one rectangle is
\[
 \mathscr X_{N,M}
 :=\prod_{n=0}^{N}L^2(0,T_*;H^{M+1})
 \times\prod_{n=0}^{N}L^2(0,T_*;H^{M-1})
 \times\prod_{n=0}^{N}L^1(0,T_*;H^M).
\]
The restriction maps make these spaces an inverse system, with
\[
\begin{aligned}
 \mathscr X_\infty
 :=\varprojlim_{N,M}\mathscr X_{N,M}
 =\bigl\{(Z_{N,M})_{N,M}:{}&
 Z_{N,M}\in\mathscr X_{N,M},\\[-2pt]
 &\rho_{N',M'}^{N,M}Z_{N,M}=Z_{N',M'}
 \ \text{if }N'\le N,\ M'\le M\bigr\}.
\end{aligned}
\]
The datum-specific intersection is the compatible family
\[
 \mathcal I[u_0]
 :=\bigl(\mathcal R_{N,M}[u_0;T_*]\bigr)_{N,M\in\mathbb N_0}
 \in\mathscr X_\infty,
 \qquad
 u\in\bigcap_{r,m<\infty}H^r(0,T_*;H^m(\mathbb R^3)).
\]
There is no all-order estimate hidden in this notation.  Each requested
coordinate is obtained at a finite stage, and compatibility makes all of
those finite coordinates parts of one history.

Take the zeroth temporal row after this intersection has been formed.  For
every finite \(m\),
\[
 u\in L^2(0,T_*;H^{m+1}),
 \qquad
 u_t\in L^2(0,T_*;H^{m-1}),
 \qquad
 u\in C([0,T_*];H^m).
\]
The endpoint values at different heights are compatible traces of one field.
They produce a common endpoint
\[
 u_*:=u(T_*)\in\bigcap_{m<\infty}H^m(\mathbb R^3)
 =H^\infty(\mathbb R^3)
 \hookrightarrow C^\infty(\mathbb R^3).
\]
At \(T_*\), set \(V_0^*:=u_*\).  The pressure formula and binomial recurrence
then define \(Q_r^*\) and \(V_{r+1}^*\) through every finite order.  Smooth
local existence restarts the flow from \(u_*\).  Uniqueness glues that
continuation to the history already reaching \(T_*\), contradicting finite
maximality.
Hence
\[
 T_*=\infty.
\]

Only after \(\mathcal I[u_0]\) is present do the familiar finite estimates
become available as coordinate projections.  Fix any finite \(T\), a temporal
order \(r\), a spatial derivative order \(k\), and \(2\le p\le\infty\), then
choose
\[
 m>k+\frac32-\frac3p.
\]
The corresponding row gives
\[
 V_r\in C([0,T];H^m)
 \hookrightarrow L^\infty(0,T;W^{k,p}).
\]
For instance, the \(m=3\) coordinate yields
\[
 \int_0^{T}\|\nabla u(t)\|_\infty\,dt
 \le T\|u\|_{L^\infty(0,T;W^{1,\infty})}<\infty.
\]
This estimate appears because the fluid already occupies the compatible
intersection; it is a finite readout of that object, not the source of it.
Its source remains visible in the full dependence
\[
\begin{gathered}
 u_0
 \longrightarrow
 \underbrace{\{(V_r,Q_r)\}_{r\ge0}}_{\text{pressure-complete mixed jet}}
 \longrightarrow
 \underbrace{\left\{
  \begin{gathered}
   V_n\in L^2(0,T_*;H^{M+1}),\\[-1pt]
   V_{n+1}\in L^2(0,T_*;H^{M-1})
  \end{gathered}
 \right\}_{\substack{0\le n\le N\\N,M<\infty}}}_{
 \text{whole-terminal adjacent rows}}
 \\
 \longrightarrow
 \underbrace{\{\mathcal R_{N,M}[u_0;T_*]\}_{N,M<\infty}}_{
 \text{finite compatible rectangles}}
 \longrightarrow
 \underbrace{\mathcal I[u_0]}_{\text{projective intersection}}
 \longrightarrow
 \underbrace{\left(
  u_*\in H^\infty,
  \{(V_r^*,Q_r^*)\}_{r\ge0}
 \right)}_{\substack{\text{common endpoint and pressure-complete}\\
                     \text{endpoint jet}}}
 \\
 \longrightarrow \text{restart and uniqueness}
 \longrightarrow T_*=\infty.
\end{gathered}
\]

\endgroup
